
\begin{figure*}[t]
  \centering
  \begin{subfigure}{0.35\textwidth}
    \centering
    \includegraphics[width=\linewidth]{Figs/DES.drawio.pdf}
    \caption{Normal DES}
    \label{fig:des}
  \end{subfigure}%
  \hspace{2em}
  \begin{subfigure}{0.35\textwidth}
    \centering
    \includegraphics[width=\linewidth]{Figs/DES2.drawio.pdf}
    \caption{DES running real applications}
    \label{fig:des-shadow}
  \end{subfigure}
  \vspace{-0.5em}
  \caption{Comparison of conventional discrete-event simulation (DES) and DES with real applications. (a) Normal DES: a network simulator translates network traffic and an abstract node model into discrete events scheduled on a unified timeline. (b) DES running real applications: the simulator retains the same event-driven network core, but replaces the abstract node model with a lightweight virtualization layer that executes real client applications. Application and network events coexist on a single time axis under a shared simulation clock.}
  \label{fig:DESs}
\end{figure*}

\section{Event-Driven Testnet}
\label{sec:edt}

This section presents the key idea behind the \emph{Event-Driven Testnet} (EDT), a hybrid \emph{testbed} that executes real blockchain client software within an event-driven network simulator. EDT treats protocol time as a first-class input controlled by the test runner, rather than as an uncontrollable background parameter bound to wall-clock execution. The design combines the performance advantages of simulators with the implementation fidelity of traditional testnets.

\subsection{Key Idea: Best-of-Both}
\label{subsec:edt-keyidea}

To address the inefficiency of testing \emph{strongly time-controlled} consensus protocols, EDT combines the strengths of conventional simulators and testnets while mitigating their respective limitations.

\noindent\textbf{Simulators:} Simulators often achieve superior performance when modeling strongly time-controlled protocols. This efficiency arises from two fundamental factors.
First, simulators deliberately abstract and simplify protocol logic, focusing on essential behavioral properties rather than faithfully executing every implementation detail. This selective abstraction reduces computational cost.
Second, and more decisively, the dominant \emph{event-driven} architecture (Figure~\ref{fig:DESs}(a)) is well-suited to protocols whose state machine advances at discrete time boundaries and is otherwise idle. Event-driven execution naturally skips empty intervals, enabling dramatic acceleration of simulated time. Simulators also typically provide strong controllability over time and network schedules, which facilitates deterministic and reproducible experiments.

Despite these advantages, simulators have significant limitations: they operate on simplified logic, may diverge from production implementations over time, and thus cannot reliably surface implementation-specific bugs or capture emergent behaviors of deployed client software.

\noindent\textbf{Testnets:} In contrast, testnets execute unmodified, production-grade client software, providing high implementation fidelity. This fidelity makes testnets valuable for uncovering implementation bugs, validating experimental protocol changes, and observing genuine network dynamics under realistic configurations.

However, for strongly time-controlled protocols, testnets inherit substantial temporal redundancy. As shown by our empirical study of Ethereum PoS (Section~\ref{sec:motivation}), nodes spend most wall-clock time idle or waiting between sparse bursts of core-logic execution, which severely limits how quickly experiments can progress over many epochs.

\textbf{The EDT synthesis.} The key insight is to execute \emph{real client applications}, like what testnets do, while leveraging the event-driven control of protocol time, like what simulators do (Figure~\ref{fig:DESs}(b)). Concretely, EDT replaces long idle waiting periods between scheduled time boundaries with instantaneous event-driven time advancement, increasing the density at which the consensus state machine is exercised without sacrificing the realism of production software. In short, EDT aims to realize a best-of-both-worlds testbed: simulator-like efficiency with testnet-like fidelity.

\subsection{Design Principles}
\label{sec:design_principles}

We define design principles to systematically handle strongly time-controlled testnets with both efficiency and fidelity. EDT is only useful if it improves testing efficiency \emph{without} compromising the causal semantics of the consensus protocol or the fidelity of client behavior.
We therefore structure Event-Driven Testnet around the following design principles.

\smallskip
\noindent\textbf{P1: Semantic equivalence of protocol time.}
The event-driven execution must preserve the protocol's causal semantics induced by its logical clock (slots/epochs) and time-dependent rules.
In particular, advancing logical time must not skip or reorder mandatory time-triggered transitions, and must maintain the intended happens-before relationships between time ticks and external events (e.g., block/attestation processing, fork-choice updates).
This principle explicitly distinguishes our approach from modifying protocol parameters (e.g., shrinking the slot duration), which changes the protocol itself rather than the execution substrate.

\smallskip
\noindent\textbf{P2: Controllability for systematic testing.}
The framework should expose control knobs over logical-time progression and network scheduling, enabling systematic exploration of ``what-if'' scenarios.
Examples include pausing and resuming logical time, stepping through critical boundaries (e.g., slot/epoch ticks), and coordinating with fault-injection tools (e.g., fuzzers) under a controlled schedule.
Controllability is essential for debugging and targeted testing within specific time intervals.

\smallskip
\noindent\textbf{P3: Reproducibility by deterministic event scheduling.}
Given a fixed configuration (client launching commandlines, genesis, network topology, delay models) and a fixed random seed, the event-driven testnet should produce a deterministic event schedule and thus a reproducible execution.
Reproducibility enables precise regression debugging, which is notoriously challenging in real-time distributed testnets.
Importantly, determinism also enables a \emph{lightweight replay} mechanism that re-executes the system under the same event schedule and a specified time window, avoiding heavyweight global snapshot/checkpoint mechanisms.

\smallskip
\noindent\textbf{P4: Implementation fidelity via real clients.}
To surface implementation-specific bugs, the framework should run real, unmodified, up-to-date client binaries end-to-end (execution layer, consensus layer, and validator software when applicable).
The event-driven substrate should therefore provide sufficiently faithful OS and network abstractions to support production clients, rather than relying on simplified protocol models.

\smallskip
These principles jointly define the target envelope for EDT: compress idle wall-clock time while preserving protocol causality and client-visible behavior, and enable deterministic, controllable experimentation for strongly time-controlled consensus protocols.
  



\subsection{Design and Implementation}
\label{subsec:edt-impl}

EDT follows the design principles (P1--P4) in Section~\ref{sec:design_principles} by turning \emph{protocol time} into a first-class, controllable input while preserving the causal semantics that real clients expect.
Our system decomposes into three layers: (i) \emph{event-driven execution of real clients} using Shadow (meeting P1 and P4), (ii) \emph{a scalable network model} implemented by DVN with runahead-style parallelism, and (iii) \emph{a time-control framework} that exposes logical time as an explicit control surface (our contribution; meeting P2 and strengthening P3).
Conceptually, EDT provides a single \emph{simulation clock} shared by the network model and all client-visible time interfaces; the simulator advances this clock by executing only causally-ready work, rather than waiting for wall-clock time to pass.

\smallskip
\noindent\textbf{Shadow as the substrate (P1/P4): decoupling time while running real binaries.}
Shadow provides the key substrate that simultaneously (i) runs \emph{unmodified production binaries} and (ii) decouples their perceived time from wall-clock time.
It virtualizes each node as a lightweight host that executes real processes while intercepting their OS interactions through a shim layer.
Critically for EDT, Shadow implements \emph{time virtualization}: all client-visible time sources and time-blocking operations are redirected to the simulation clock, so the client ``sees'' logical time rather than wall-clock time.
Likewise, blocking I/O and waiting are converted into \emph{discrete wakeup conditions} scheduled on the same event timeline.
This design meets P1 by ensuring that time advances only through explicit simulator-controlled transitions on the unified timeline, and meets P4 by preserving full implementation fidelity of real clients.

\smallskip
\noindent\textbf{Network scheduling and parallelism: DVN with runahead windows.}
EDT instantiates the network model using a Distributed Virtual Network (DVN) simulator, which provides configurable topology, latency distributions, loss/adversary models, and produces a deterministic schedule of message deliveries under a fixed seed.
To scale to many hosts, DVN employs a \emph{runahead}-style parallel simulation strategy organized around a moving \emph{window} (a safe synchronization horizon).
Within a window \([t, t+\Delta)\), each host (or partition) can advance locally and process events as long as DVN can conservatively guarantee that no cross-partition delivery can arrive earlier than the window end.
When the window boundary is reached, partitions synchronize: they exchange queued cross-partition messages, advance the global frontier, and open the next window.
This yields parallelism by exploiting causally-independent event sets, and it provides a natural notion of ``safe boundary'' that we later reuse for time control and debugging.
The parallelism is complementary to Shadow's time virtualization: Shadow compresses idle time \emph{within} each host, while DVN accelerates the network core by running many hosts concurrently up to the runahead window horizon.

\smallskip
\noindent\textbf{A time-control framework (P2/P3).}
Building on Shadow (P1/P4) and DVN runahead windows, EDT exposes a \emph{time-control surface} that lets the test runner drive protocol time explicitly.
\emph{Even without any special control}, EDT already achieves significant speedups: strongly time-controlled clients spend much wall-clock time in sleeps, timer waits, and network blocking between sparse bursts of protocol logic.
In EDT these idle intervals are represented as pending wakeup conditions on the event timeline, and the simulator advances the simulation clock directly to the next causally-ready event time.
Hence, long periods of ``nothing happens'' become instantaneous time jumps, and computation is concentrated on the discrete moments when the consensus state machine actually performs work (e.g., slot/epoch transitions, message arrivals, timer expirations).

\smallskip
\noindent\emph{GDB-like control style.}
We then make logical time \emph{explicitly controllable} by exposing a small set of commands whose usage mirrors interactive debuggers (e.g., \texttt{continue}, \texttt{next}, \texttt{step}, \texttt{restart}).
For example, our implementation supports (i) pausing and resuming time progression, (ii) running until a target logical time \(t\), and (iii) stepping by one runahead \emph{window} at a time (``next-window'').
The step granularity is chosen deliberately: DVN's runahead window is the minimal unit at which hosts are guaranteed not to causally affect each other via cross-host deliveries.
Therefore, pausing at a window boundary yields a consistent global cut of the distributed execution, and stepping one window preserves determinism while keeping synchronization overhead low.

\smallskip
\noindent\emph{Distributed debugging and fuzzing potential.}
Pausing at a window boundary also enables a practical ``distributed debugging'' workflow.
At the boundary, the runtime can enumerate which hosts are scheduled to run in the next window and report the corresponding native process identifiers (PIDs) of the real clients.
A tester can then attach a debugger (e.g., \texttt{gdb -p <pid>}) to one or more target processes, inspect per-node local state and execution, and resume the simulation deterministically.
Because hosts are independent within a window under the runahead guarantee, a tester can single-step or instrument selected nodes without perturbing the causal schedule of other nodes in that window.
Similarly, a fuzzer can treat time boundaries as explicit intervention points, inject faults or schedule perturbations at controlled windows, and replay the same failing window under additional instrumentation.
